\documentclass[aspectratio=169, 11pt]{beamer}
% ============================================================================
% Preambulo compartido para presentaciones Beamer
% Estadistica 2026-II - Pontificia Universidad Javeriana
% Incluir con: % ============================================================================
% Preambulo compartido para presentaciones Beamer
% Estadistica 2026-II - Pontificia Universidad Javeriana
% Incluir con: % ============================================================================
% Preambulo compartido para presentaciones Beamer
% Estadistica 2026-II - Pontificia Universidad Javeriana
% Incluir con: \input{../preambulo_beamer.tex} despues de \documentclass
% ============================================================================

\usepackage[T1]{fontenc}
\usepackage[utf8]{inputenc}
\usepackage[spanish]{babel}
\usepackage{amsmath, amssymb, amsfonts}
\usepackage{booktabs}
\usepackage{graphicx}
\usepackage{tikz}
\usetikzlibrary{shapes.geometric, positioning, arrows.meta, calc}
\usepackage{pgfplots}
\pgfplotsset{compat=1.18}
\usepackage{listings}
\usepackage{xcolor}
\usepackage{hyperref}
\usepackage{multicol}

% --- Tema Beamer (minimalista) ---
\usetheme{default}
\usecolortheme{default}
\usefonttheme{default}
\setbeamertemplate{navigation symbols}{}
\setbeamertemplate{caption}[numbered]
\setbeamertemplate{footline}{%
  \hspace{1em}{\tiny\url{https://github.com/polanco-jaime/Curso-Estadistica-2026-3}}\hfill\usebeamerfont{footline}\insertframenumber/\inserttotalframenumber\hspace{1em}\vspace{0.5em}%
}
\setbeamertemplate{itemize items}[circle]
\setbeamertemplate{enumerate items}[default]

% --- Colores institucionales (sutiles) ---
\definecolor{javeriana}{RGB}{0, 62, 105}
\definecolor{javerianalight}{RGB}{0, 105, 170}
\definecolor{codigobg}{RGB}{245, 245, 245}
\definecolor{codigofg}{RGB}{40, 40, 40}
\definecolor{comentario}{RGB}{100, 100, 100}
\definecolor{keyword}{RGB}{0, 100, 150}
\definecolor{string}{RGB}{180, 60, 30}

\setbeamercolor{title}{fg=javeriana}
\setbeamercolor{frametitle}{fg=javeriana}
\setbeamercolor{structure}{fg=javeriana}
\setbeamercolor{block title}{fg=white, bg=javeriana!75}
\setbeamercolor{block body}{bg=javeriana!5}
\setbeamercolor{block title alerted}{fg=white, bg=red!70!black}
\setbeamercolor{block body alerted}{bg=red!5}
\setbeamercolor{block title example}{fg=white, bg=green!40!black}
\setbeamercolor{block body example}{bg=green!5}
\setbeamercolor{item}{fg=javeriana}

\setbeamerfont{title}{series=\bfseries, size=\Large}
\setbeamerfont{frametitle}{series=\bfseries}

% --- Configuracion de listings para R ---
\lstset{
  language=R,
  basicstyle=\ttfamily\footnotesize,
  backgroundcolor=\color{codigobg},
  commentstyle=\color{comentario}\itshape,
  keywordstyle=\color{keyword}\bfseries,
  stringstyle=\color{string},
  numbers=none,
  frame=single,
  rulecolor=\color{javeriana!30},
  breaklines=true,
  showstringspaces=false,
  columns=flexible,
  morekeywords={library, ggplot, aes, geom_histogram, geom_boxplot,
    geom_point, geom_smooth, geom_density, geom_bar, geom_line,
    geom_vline, geom_hline, geom_errorbar, labs, theme_minimal,
    facet_wrap, scale_fill_manual, scale_color_manual,
    summarise, group_by, filter, mutate, select, arrange,
    read.csv, head, str, summary, mean, median, sd, var, cor,
    lm, t.test, chisq.test, wilcox.test, kruskal.test, aov,
    pnorm, qnorm, dnorm, rnorm, qt, pt, qchisq, pchisq, qf, pf,
    confint, coeftest, vcovHC, bptest, vif, cooks.distance,
    TukeyHSD, table, prop.table}
}

% --- Comandos personalizados ---
\newcommand{\R}{\texttt{R}}
\newcommand{\code}[1]{\texttt{\footnotesize #1}}
\newcommand{\variable}[1]{\texttt{#1}}
\newcommand{\paquete}[1]{\texttt{#1}}

% --- Informacion del curso ---
\institute{Pontificia Universidad Javeriana\\
  Facultad de Ciencias Econ\'omicas y Administrativas}
\date{2026-II}
 despues de \documentclass
% ============================================================================

\usepackage[T1]{fontenc}
\usepackage[utf8]{inputenc}
\usepackage[spanish]{babel}
\usepackage{amsmath, amssymb, amsfonts}
\usepackage{booktabs}
\usepackage{graphicx}
\usepackage{tikz}
\usetikzlibrary{shapes.geometric, positioning, arrows.meta, calc}
\usepackage{pgfplots}
\pgfplotsset{compat=1.18}
\usepackage{listings}
\usepackage{xcolor}
\usepackage{hyperref}
\usepackage{multicol}

% --- Tema Beamer (minimalista) ---
\usetheme{default}
\usecolortheme{default}
\usefonttheme{default}
\setbeamertemplate{navigation symbols}{}
\setbeamertemplate{caption}[numbered]
\setbeamertemplate{footline}{%
  \hspace{1em}{\tiny\url{https://github.com/polanco-jaime/Curso-Estadistica-2026-3}}\hfill\usebeamerfont{footline}\insertframenumber/\inserttotalframenumber\hspace{1em}\vspace{0.5em}%
}
\setbeamertemplate{itemize items}[circle]
\setbeamertemplate{enumerate items}[default]

% --- Colores institucionales (sutiles) ---
\definecolor{javeriana}{RGB}{0, 62, 105}
\definecolor{javerianalight}{RGB}{0, 105, 170}
\definecolor{codigobg}{RGB}{245, 245, 245}
\definecolor{codigofg}{RGB}{40, 40, 40}
\definecolor{comentario}{RGB}{100, 100, 100}
\definecolor{keyword}{RGB}{0, 100, 150}
\definecolor{string}{RGB}{180, 60, 30}

\setbeamercolor{title}{fg=javeriana}
\setbeamercolor{frametitle}{fg=javeriana}
\setbeamercolor{structure}{fg=javeriana}
\setbeamercolor{block title}{fg=white, bg=javeriana!75}
\setbeamercolor{block body}{bg=javeriana!5}
\setbeamercolor{block title alerted}{fg=white, bg=red!70!black}
\setbeamercolor{block body alerted}{bg=red!5}
\setbeamercolor{block title example}{fg=white, bg=green!40!black}
\setbeamercolor{block body example}{bg=green!5}
\setbeamercolor{item}{fg=javeriana}

\setbeamerfont{title}{series=\bfseries, size=\Large}
\setbeamerfont{frametitle}{series=\bfseries}

% --- Configuracion de listings para R ---
\lstset{
  language=R,
  basicstyle=\ttfamily\footnotesize,
  backgroundcolor=\color{codigobg},
  commentstyle=\color{comentario}\itshape,
  keywordstyle=\color{keyword}\bfseries,
  stringstyle=\color{string},
  numbers=none,
  frame=single,
  rulecolor=\color{javeriana!30},
  breaklines=true,
  showstringspaces=false,
  columns=flexible,
  morekeywords={library, ggplot, aes, geom_histogram, geom_boxplot,
    geom_point, geom_smooth, geom_density, geom_bar, geom_line,
    geom_vline, geom_hline, geom_errorbar, labs, theme_minimal,
    facet_wrap, scale_fill_manual, scale_color_manual,
    summarise, group_by, filter, mutate, select, arrange,
    read.csv, head, str, summary, mean, median, sd, var, cor,
    lm, t.test, chisq.test, wilcox.test, kruskal.test, aov,
    pnorm, qnorm, dnorm, rnorm, qt, pt, qchisq, pchisq, qf, pf,
    confint, coeftest, vcovHC, bptest, vif, cooks.distance,
    TukeyHSD, table, prop.table}
}

% --- Comandos personalizados ---
\newcommand{\R}{\texttt{R}}
\newcommand{\code}[1]{\texttt{\footnotesize #1}}
\newcommand{\variable}[1]{\texttt{#1}}
\newcommand{\paquete}[1]{\texttt{#1}}

% --- Informacion del curso ---
\institute{Pontificia Universidad Javeriana\\
  Facultad de Ciencias Econ\'omicas y Administrativas}
\date{2026-II}
 despues de \documentclass
% ============================================================================

\usepackage[T1]{fontenc}
\usepackage[utf8]{inputenc}
\usepackage[spanish]{babel}
\usepackage{amsmath, amssymb, amsfonts}
\usepackage{booktabs}
\usepackage{graphicx}
\usepackage{tikz}
\usetikzlibrary{shapes.geometric, positioning, arrows.meta, calc}
\usepackage{pgfplots}
\pgfplotsset{compat=1.18}
\usepackage{listings}
\usepackage{xcolor}
\usepackage{hyperref}
\usepackage{multicol}

% --- Tema Beamer (minimalista) ---
\usetheme{default}
\usecolortheme{default}
\usefonttheme{default}
\setbeamertemplate{navigation symbols}{}
\setbeamertemplate{caption}[numbered]
\setbeamertemplate{footline}{%
  \hspace{1em}{\tiny\url{https://github.com/polanco-jaime/Curso-Estadistica-2026-3}}\hfill\usebeamerfont{footline}\insertframenumber/\inserttotalframenumber\hspace{1em}\vspace{0.5em}%
}
\setbeamertemplate{itemize items}[circle]
\setbeamertemplate{enumerate items}[default]

% --- Colores institucionales (sutiles) ---
\definecolor{javeriana}{RGB}{0, 62, 105}
\definecolor{javerianalight}{RGB}{0, 105, 170}
\definecolor{codigobg}{RGB}{245, 245, 245}
\definecolor{codigofg}{RGB}{40, 40, 40}
\definecolor{comentario}{RGB}{100, 100, 100}
\definecolor{keyword}{RGB}{0, 100, 150}
\definecolor{string}{RGB}{180, 60, 30}

\setbeamercolor{title}{fg=javeriana}
\setbeamercolor{frametitle}{fg=javeriana}
\setbeamercolor{structure}{fg=javeriana}
\setbeamercolor{block title}{fg=white, bg=javeriana!75}
\setbeamercolor{block body}{bg=javeriana!5}
\setbeamercolor{block title alerted}{fg=white, bg=red!70!black}
\setbeamercolor{block body alerted}{bg=red!5}
\setbeamercolor{block title example}{fg=white, bg=green!40!black}
\setbeamercolor{block body example}{bg=green!5}
\setbeamercolor{item}{fg=javeriana}

\setbeamerfont{title}{series=\bfseries, size=\Large}
\setbeamerfont{frametitle}{series=\bfseries}

% --- Configuracion de listings para R ---
\lstset{
  language=R,
  basicstyle=\ttfamily\footnotesize,
  backgroundcolor=\color{codigobg},
  commentstyle=\color{comentario}\itshape,
  keywordstyle=\color{keyword}\bfseries,
  stringstyle=\color{string},
  numbers=none,
  frame=single,
  rulecolor=\color{javeriana!30},
  breaklines=true,
  showstringspaces=false,
  columns=flexible,
  morekeywords={library, ggplot, aes, geom_histogram, geom_boxplot,
    geom_point, geom_smooth, geom_density, geom_bar, geom_line,
    geom_vline, geom_hline, geom_errorbar, labs, theme_minimal,
    facet_wrap, scale_fill_manual, scale_color_manual,
    summarise, group_by, filter, mutate, select, arrange,
    read.csv, head, str, summary, mean, median, sd, var, cor,
    lm, t.test, chisq.test, wilcox.test, kruskal.test, aov,
    pnorm, qnorm, dnorm, rnorm, qt, pt, qchisq, pchisq, qf, pf,
    confint, coeftest, vcovHC, bptest, vif, cooks.distance,
    TukeyHSD, table, prop.table}
}

% --- Comandos personalizados ---
\newcommand{\R}{\texttt{R}}
\newcommand{\code}[1]{\texttt{\footnotesize #1}}
\newcommand{\variable}[1]{\texttt{#1}}
\newcommand{\paquete}[1]{\texttt{#1}}

% --- Informacion del curso ---
\institute{Pontificia Universidad Javeriana\\
  Facultad de Ciencias Econ\'omicas y Administrativas}
\date{2026-II}

\title{Sesión 11: Diagnósticos de Regresión}
\subtitle{Validando los supuestos del modelo OLS}
\author{Prof.\ Jaime Polanco Jim\'enez}
\date{Viernes 6 de noviembre de 2026}

\begin{document}

\begin{frame}
\titlepage
\end{frame}

\begin{frame}{Agenda}
\tableofcontents
\end{frame}

% ============================================================================
\section{¿Por qué importan los diagnósticos?}
% ============================================================================

\begin{frame}{Los supuestos del modelo OLS}
\begin{block}{Recapitulación: el modelo clásico de regresión lineal}
$$Y_i = \beta_0 + \beta_1 X_{1i} + \beta_2 X_{2i} + \cdots + \beta_k X_{ki} + \varepsilon_i$$
\end{block}

\textbf{Supuestos clave (Gauss-Markov):}
\begin{enumerate}
\item \textbf{Linealidad}: la relación entre $Y$ y las $X$ es lineal
\item \textbf{Exogeneidad}: $E(\varepsilon_i | X) = 0$
\item \textbf{Homocedasticidad}: $\text{Var}(\varepsilon_i | X) = \sigma^2$ (constante)
\item \textbf{No autocorrelación}: $\text{Cor}(\varepsilon_i, \varepsilon_j) = 0$ para $i \neq j$
\item \textbf{Normalidad} (para inferencia): $\varepsilon_i \sim N(0, \sigma^2)$
\end{enumerate}

\vspace{0.3cm}
\begin{alertblock}{Resultado fundamental}
Si se cumplen los supuestos 1-4, OLS es \textbf{BLUE} (Best Linear Unbiased Estimator).
Si además se cumple 5, las pruebas $t$ y $F$ son válidas.
\end{alertblock}
\end{frame}

\begin{frame}{¿Qué pasa si se violan los supuestos?}
\small
\begin{center}
\renewcommand{\arraystretch}{1.2}
\begin{tabular}{lll}
\toprule
\textbf{Supuesto violado} & \textbf{Consecuencia} & \textbf{Síntoma} \\
\midrule
Linealidad & Estimadores sesgados & Patr\'on en residuos \\
Exogeneidad & Estimadores sesgados & Variable omitida \\
Homocedasticidad & EE incorrectos, ineficiente & Embudo en residuos \\
No autocorrelaci\'on & EE incorrectos & Patr\'on temporal \\
Normalidad & IC y pruebas inv\'alidas & Q-Q desviado \\
\bottomrule
\end{tabular}
\end{center}

\begin{exampleblock}{Ejemplo: heterocedasticidad}
\small
Los $\hat{\beta}$ siguen insesgados, pero los errores est\'andar son \textbf{incorrectos}
$\Rightarrow$ los IC y las pruebas $t$ son \textbf{inv\'alidas}.

\smallskip
Piensen en las notas de ingl\'es: la variabilidad es mayor entre estudiantes que trabajan medio tiempo --- sus resultados son m\'as impredecibles que los de quienes se dedican solo a estudiar. Eso es heterocedasticidad.
\end{exampleblock}
\end{frame}

\begin{frame}{El papel de los diagnósticos}
\begin{block}{Diagnósticos = verificar supuestos \textit{después} de estimar}
Los diagnósticos nos permiten detectar problemas con el modelo ajustado y decidir si necesitamos:
\begin{itemize}
\item Transformar variables
\item Agregar/eliminar variables
\item Usar errores estándar robustos
\item Cambiar de modelo (ej: pasar a logit, poisson, etc.)
\end{itemize}
\end{block}

\vspace{0.15cm}
\begin{center}
\begin{tikzpicture}[
  node distance=1.2cm,
  box/.style={rectangle, draw=javeriana, fill=javeriana!10, thick, minimum width=2.2cm, minimum height=0.7cm, rounded corners, font=\small}
]
\node[box] (est) {Estimar modelo};
\node[box, right=of est] (diag) {Diagnosticar};
\node[box, right=of diag] (corr) {Corregir};
\node[box, right=of corr] (reest) {Re-estimar};

\draw[->, thick, javeriana] (est) -- (diag);
\draw[->, thick, javeriana] (diag) -- (corr);
\draw[->, thick, javeriana] (corr) -- (reest);
\draw[->, thick, javeriana] (reest) to[bend right=40] (diag);
\end{tikzpicture}
\end{center}
\end{frame}

\begin{frame}{Herramientas diagnósticas en R}
\textbf{Principales funciones:}
\begin{itemize}
\item \code{plot(modelo)}: genera 4 gráficos diagnósticos automáticamente
\item \code{residuals(modelo)}: extrae los residuos $e_i = Y_i - \hat{Y}_i$
\item \code{fitted(modelo)}: extrae los valores ajustados $\hat{Y}_i$
\item \code{rstandard(modelo)}: residuos estandarizados
\item \code{cooks.distance(modelo)}: distancia de Cook para cada observación
\item \code{lmtest::bptest(modelo)}: test de Breusch-Pagan (heterocedasticidad)
\item \code{lmtest::dwtest(modelo)}: test de Durbin-Watson (autocorrelación)
\item \code{car::vif(modelo)}: factor de inflación de la varianza (multicolinealidad)
\end{itemize}

\vspace{0.3cm}
\begin{alertblock}{Estrategia}
Comenzar con \code{plot(modelo)} y luego profundizar con tests formales según lo que se observe.
\end{alertblock}
\end{frame}

% ============================================================================
\section{Los 4 gráficos diagnósticos de R}
% ============================================================================

\begin{frame}[fragile]{Los 4 gráficos de \code{plot(modelo)}}
\begin{block}{Generación automática en R}
\begin{lstlisting}
# Ajustar modelo
modelo <- lm(MOD_INGLES_PUNT ~ ESTU_ESTRATO + PRIVADA +
             ESTU_GENERO + PUNT_INGLES_S11, data = ni_data)

# Generar los 4 graficos diagnosticos
par(mfrow = c(2, 2))
plot(modelo)
\end{lstlisting}
\end{block}

\textbf{Los 4 gráficos son:}
\begin{enumerate}
\item \textbf{Residuals vs Fitted}: detecta no linealidad y heterocedasticidad
\item \textbf{Normal Q-Q}: verifica normalidad de residuos
\item \textbf{Scale-Location}: detecta heterocedasticidad
\item \textbf{Residuals vs Leverage}: identifica observaciones influyentes
\end{enumerate}
\end{frame}

\begin{frame}{Gráfico 1: Residuals vs Fitted}
\textbf{¿Qué muestra?}
\begin{itemize}
\item Eje X: valores ajustados $\hat{Y}_i$
\item Eje Y: residuos $e_i$
\item Línea roja: tendencia suavizada (loess)
\end{itemize}

\vspace{0.3cm}
\textbf{¿Qué buscar?}
\begin{itemize}
\item \textcolor{javerianalight}{Ideal}: puntos dispersos aleatoriamente alrededor de $y=0$, línea roja horizontal
\item \textcolor{red}{Problema 1}: línea roja con curvatura $\Rightarrow$ no linealidad
\item \textcolor{red}{Problema 2}: forma de embudo $\Rightarrow$ heterocedasticidad
\end{itemize}
\end{frame}

\begin{frame}{Gráfico 1: Patrones visuales}
\begin{columns}
\begin{column}{0.48\textwidth}
\begin{exampleblock}{Patrón ideal}
\begin{center}
\begin{tikzpicture}[scale=0.7]
\begin{axis}[
  width=5cm, height=4cm,
  xlabel={Fitted values},
  ylabel={Residuals},
  grid=major,
  xmin=120, xmax=180,
  ymin=-30, ymax=30
]
\addplot[only marks, mark=*, mark size=1pt, opacity=0.5]
  coordinates {(150,5) (145,-10) (155,15) (140,-5) (160,10) (148,-15) (152,8) (158,-12) (142,20)};
\addplot[red, thick, domain=120:180] {0};
\end{axis}
\end{tikzpicture}
\end{center}
Dispersión constante alrededor de cero.
\end{exampleblock}
\end{column}

\begin{column}{0.48\textwidth}
\begin{alertblock}{Patrón problemático}
\begin{center}
\begin{tikzpicture}[scale=0.7]
\begin{axis}[
  width=5cm, height=4cm,
  xlabel={Fitted values},
  ylabel={Residuals},
  grid=major,
  xmin=120, xmax=180,
  ymin=-30, ymax=30
]
\addplot[only marks, mark=*, mark size=1pt, opacity=0.5]
  coordinates {(130,-2) (135,5) (140,-8) (145,15) (150,-5) (155,18) (160,-10) (165,25) (170,-15)};
\addplot[red, thick, domain=120:180] {0.3*x - 45};
\end{axis}
\end{tikzpicture}
\end{center}
Curvatura en la línea roja.
\end{alertblock}
\end{column}
\end{columns}
\end{frame}

\begin{frame}{Gráfico 2: Normal Q-Q}
\textbf{¿Qué muestra?}
\begin{itemize}
\item Eje X: cuantiles teóricos de $N(0,1)$
\item Eje Y: residuos estandarizados
\item Línea diagonal: referencia para normalidad perfecta
\end{itemize}

\vspace{0.3cm}
\textbf{¿Qué buscar?}
\begin{itemize}
\item \textcolor{javerianalight}{Ideal}: puntos siguen la línea diagonal
\item \textcolor{red}{Problema}: desviación sistemática:
  \begin{itemize}
  \item Colas pesadas: puntos se alejan en los extremos
  \item Asimetría: curvatura en S
  \end{itemize}
\end{itemize}

\vspace{0.3cm}
\begin{block}{Interpretación}
Desviaciones leves son normales. Preocuparse solo si la desviación es marcada y afecta muchos puntos.
\end{block}
\end{frame}

\begin{frame}{Gráfico 2: Patrones Q-Q}
\begin{columns}
\begin{column}{0.48\textwidth}
\begin{exampleblock}{Normalidad (bueno)}
\begin{center}
\begin{tikzpicture}[scale=0.7]
\begin{axis}[
  width=5cm, height=4cm,
  xlabel={Theoretical Quantiles},
  ylabel={Standardized residuals},
  grid=major,
  xmin=-3, xmax=3,
  ymin=-3, ymax=3
]
\addplot[only marks, mark=*, mark size=1pt, opacity=0.5]
  coordinates {(-2.5,-2.4) (-2.0,-2.1) (-1.5,-1.6) (-1.0,-0.9) (-0.5,-0.6) (0,0.1) (0.5,0.4) (1.0,1.1) (1.5,1.4) (2.0,2.2) (2.5,2.6)};
\addplot[dashed, thick, domain=-3:3] {x};
\end{axis}
\end{tikzpicture}
\end{center}
Puntos cerca de la línea diagonal.
\end{exampleblock}
\end{column}

\begin{column}{0.48\textwidth}
\begin{alertblock}{Colas pesadas (malo)}
\begin{center}
\begin{tikzpicture}[scale=0.7]
\begin{axis}[
  width=5cm, height=4cm,
  xlabel={Theoretical Quantiles},
  ylabel={Standardized residuals},
  grid=major,
  xmin=-3, xmax=3,
  ymin=-4, ymax=4
]
\addplot[only marks, mark=*, mark size=1pt, opacity=0.5]
  coordinates {(-2.5,-3.5) (-2.0,-2.1) (-1.5,-1.4) (-1.0,-0.9) (-0.5,-0.4) (0,0.1) (0.5,0.5) (1.0,0.9) (1.5,1.5) (2.0,2.2) (2.5,3.6)};
\addplot[dashed, thick, domain=-3:3] {x};
\end{axis}
\end{tikzpicture}
\end{center}
Puntos se alejan en los extremos.
\end{alertblock}
\end{column}
\end{columns}
\end{frame}

\begin{frame}{Gráfico 3: Scale-Location}
\textbf{¿Qué muestra?}
\begin{itemize}
\item Eje X: valores ajustados $\hat{Y}_i$
\item Eje Y: $\sqrt{|\text{residuos estandarizados}|}$
\item Línea roja: tendencia suavizada
\end{itemize}

\vspace{0.3cm}
\textbf{¿Para qué sirve?}
\begin{itemize}
\item Detecta heterocedasticidad de manera más clara que Residuals vs Fitted
\item Al elevar al cuadrado (bueno, raíz cuadrada), se enfatiza la dispersión
\end{itemize}

\vspace{0.3cm}
\textbf{¿Qué buscar?}
\begin{itemize}
\item \textcolor{javerianalight}{Ideal}: línea roja horizontal, puntos con dispersión constante
\item \textcolor{red}{Problema}: línea roja con pendiente $\Rightarrow$ heterocedasticidad
\end{itemize}
\end{frame}

\begin{frame}{Gráfico 3: Patrones Scale-Location}
\begin{columns}
\begin{column}{0.48\textwidth}
\begin{exampleblock}{Homocedasticidad}
\begin{center}
\begin{tikzpicture}[scale=0.7]
\begin{axis}[
  width=5cm, height=3.5cm,
  xlabel={Fitted values},
  ylabel={$\sqrt{|\text{Std. residuals}|}$},
  grid=major,
  xmin=120, xmax=180,
  ymin=0, ymax=2
]
\addplot[only marks, mark=*, mark size=1pt, opacity=0.5]
  coordinates {(130,0.5) (140,1.2) (150,0.8) (160,1.0) (170,0.7) (145,1.5) (155,0.9)};
\addplot[red, thick, domain=120:180] {1.0};
\end{axis}
\end{tikzpicture}
\end{center}
Línea roja horizontal.
\end{exampleblock}
\end{column}

\begin{column}{0.48\textwidth}
\begin{alertblock}{Heterocedasticidad}
\begin{center}
\begin{tikzpicture}[scale=0.7]
\begin{axis}[
  width=5cm, height=3.5cm,
  xlabel={Fitted values},
  ylabel={$\sqrt{|\text{Std. residuals}|}$},
  grid=major,
  xmin=120, xmax=180,
  ymin=0, ymax=2.5
]
\addplot[only marks, mark=*, mark size=1pt, opacity=0.5]
  coordinates {(130,0.3) (140,0.6) (150,1.0) (160,1.4) (170,2.0) (145,0.5) (155,1.2)};
\addplot[red, thick, domain=120:180] {0.02*x - 2.1};
\end{axis}
\end{tikzpicture}
\end{center}
Línea roja con pendiente.
\end{alertblock}
\end{column}
\end{columns}
\end{frame}

\begin{frame}{Gráfico 4: Residuals vs Leverage}
\textbf{¿Qué muestra?}
\begin{itemize}
\item Eje X: leverage (apalancamiento) $h_{ii}$
\item Eje Y: residuos estandarizados
\item Líneas punteadas: isolíneas de distancia de Cook
\end{itemize}

\vspace{0.3cm}
\textbf{Conceptos clave:}
\begin{itemize}
\item \textbf{Leverage}: qué tan extremo es $X_i$ respecto al centro de los datos
\item \textbf{Residuo}: qué tan lejos está $Y_i$ de $\hat{Y}_i$
\item \textbf{Observación influyente}: alto leverage + alto residuo
\end{itemize}

\vspace{0.3cm}
\textbf{¿Qué buscar?}
\begin{itemize}
\item Puntos fuera de las líneas de Cook $D > 0.5$ o $D > 1$
\item Estas observaciones cambian mucho los coeficientes
\end{itemize}
\end{frame}

\begin{frame}{Gráfico 4: Interpretación}
\begin{center}
\begin{tikzpicture}[scale=0.9]
\begin{axis}[
  width=8cm, height=6cm,
  xlabel={Leverage},
  ylabel={Standardized residuals},
  grid=major,
  xmin=0, xmax=0.15,
  ymin=-3, ymax=3,
  legend pos=north west
]
% Puntos normales
\addplot[only marks, mark=*, mark size=1.5pt, opacity=0.6, color=javerianalight]
  coordinates {(0.02,0.5) (0.03,-0.8) (0.025,1.2) (0.04,-0.3) (0.035,0.9)};
\addlegendentry{Normal}

% Punto influyente
\addplot[only marks, mark=*, mark size=3pt, color=red]
  coordinates {(0.12,2.5)};
\addlegendentry{Influyente}

% Lineas de Cook
\addplot[dashed, domain=0:0.15, samples=50] {sqrt(0.5*5*x/(1-x))};
\addplot[dashed, domain=0:0.15, samples=50] {-sqrt(0.5*5*x/(1-x))};
\end{axis}
\end{tikzpicture}
\end{center}

\vspace{0.3cm}
El punto rojo tiene alto leverage y alto residuo $\Rightarrow$ influyente
\end{frame}

\begin{frame}[fragile]{Ejemplo en R: generando los 4 gráficos}
\begin{lstlisting}
# Cargar datos
ni_data <- read.csv("icfes_ni.csv")

# Estimar modelo
modelo <- lm(MOD_INGLES_PUNT ~ ESTU_ESTRATO + PRIVADA +
             ESTU_GENERO + PUNT_INGLES_S11, data = ni_data)

# Generar los 4 graficos diagnosticos
par(mfrow = c(2, 2))  # 2x2 grid
plot(modelo)
par(mfrow = c(1, 1))  # volver a 1x1
\end{lstlisting}
\end{frame}

\begin{frame}[fragile]{Ejemplo en R: gráficos individuales}
\begin{lstlisting}
# Para ver cada grafico individualmente:
plot(modelo, which = 1)  # Residuals vs Fitted
plot(modelo, which = 2)  # Q-Q plot
plot(modelo, which = 3)  # Scale-Location
plot(modelo, which = 5)  # Residuals vs Leverage (es el #5, no el #4)
\end{lstlisting}

\vspace{0.3cm}
\begin{alertblock}{Nota}
El gráfico 4 por defecto es Cook's distance, pero el más útil para observaciones influyentes es el 5 (Residuals vs Leverage).
\end{alertblock}
\end{frame}

\begin{frame}{Interpretación conjunta de los gráficos}
\begin{center}
\renewcommand{\arraystretch}{1.4}
\begin{tabular}{lll}
\toprule
\textbf{Gráfico} & \textbf{Detecta} & \textbf{Solución} \\
\midrule
Residuals vs Fitted & No linealidad & Transformar Y o X \\
 & Heterocedasticidad & Errores robustos \\
\midrule
Normal Q-Q & No normalidad & Transformar Y, o confiar TCL \\
\midrule
Scale-Location & Heterocedasticidad & Errores robustos \\
\midrule
Residuals vs Leverage & Obs. influyentes & Investigar, quizás eliminar \\
\bottomrule
\end{tabular}
\end{center}
\end{frame}

\begin{frame}{Estrategia práctica de diagnóstico}
\begin{exampleblock}{Pasos a seguir}
\begin{enumerate}
\item Ajustar el modelo
\item Generar los 4 gr\'aficos con \code{plot(modelo)}
\item Si se observan problemas, realizar tests formales (Breusch-Pagan, etc.)
\item Aplicar correcciones seg\'un sea necesario
\item Re-estimar y verificar de nuevo
\end{enumerate}

\smallskip
\small Es como cuando suben una foto a Instagram: primero la toman, luego revisan si sali\'o bien (diagn\'osticos), aplican filtros si es necesario (correcciones) y vuelven a revisar antes de publicar.
\end{exampleblock}
\end{frame}

% ============================================================================
\section{Heterocedasticidad y Breusch-Pagan}
% ============================================================================

\begin{frame}{El supuesto de homocedasticidad}
\begin{block}{Definici\'on formal}
\textbf{Homocedasticidad}: la varianza del error es constante para todas las observaciones.
$$\text{Var}(\varepsilon_i | X_1, X_2, \ldots, X_k) = \sigma^2 \quad \text{para todo } i$$
\end{block}

\textbf{Heterocedasticidad}: la varianza del error NO es constante.
$$\text{Var}(\varepsilon_i | X) = \sigma_i^2 \quad \text{(var\'ia con } i \text{)}$$
\end{frame}

\begin{frame}{El supuesto de homocedasticidad (cont.)}
\begin{columns}
\begin{column}{0.48\textwidth}
\begin{exampleblock}{Ejemplo de homocedasticidad}
Estudiantes de todos los estratos tienen la misma dispersi\'on en ingl\'es despu\'es de controlar por las $X$. Como si todos tuvieran el mismo nivel de incertidumbre en sus notas.
\end{exampleblock}
\end{column}

\begin{column}{0.48\textwidth}
\begin{alertblock}{Ejemplo de heterocedasticidad}
Estudiantes de estratos bajos tienen mayor variabilidad en ingl\'es que los de estratos altos (despu\'es de controlar). Es como el screen time: entre quienes pasan m\'as de 6 horas en redes, la variabilidad de notas es enorme --- algunos rinden bien y otros no.
\end{alertblock}
\end{column}
\end{columns}
\end{frame}

\begin{frame}{Consecuencias de la heterocedasticidad}
\begin{alertblock}{Impacto en OLS}
\begin{itemize}
\item \textcolor{javerianalight}{Los coeficientes $\hat{\beta}$ siguen siendo insesgados}
\item \textcolor{red}{Pero ya NO son eficientes (BLUE se pierde)}
\item \textcolor{red}{Los errores estándar calculados por OLS son incorrectos}
\item \textcolor{red}{Los IC y las pruebas $t$ y $F$ son inválidas}
\end{itemize}
\end{alertblock}

\vspace{0.15cm}
\begin{block}{¿Por qué los errores estándar son incorrectos?}
La fórmula estándar asume varianza constante $\sigma^2$: $\text{Var}(\hat{\beta}) = \sigma^2 (X'X)^{-1}$.

Con heterocedasticidad, la varianza verdadera es:
$$\text{Var}(\hat{\beta}) = (X'X)^{-1} X' \Omega X (X'X)^{-1}$$
donde $\Omega = \text{diag}(\sigma_1^2, \sigma_2^2, \ldots, \sigma_n^2)$.
\end{block}
\end{frame}

\begin{frame}{Test de Breusch-Pagan}
\begin{block}{Hipótesis}
$$H_0: \text{Homocedasticidad (varianza constante)} \qquad H_1: \text{Heterocedasticidad}$$
\end{block}

\textbf{Procedimiento:}
\begin{enumerate}
\item Estimar el modelo original y obtener los residuos $e_i$
\item Calcular $e_i^2$ (residuos al cuadrado)
\item Regresar $e_i^2$ contra TODAS las variables $X$ del modelo original:
  $$e_i^2 = \gamma_0 + \gamma_1 X_{1i} + \gamma_2 X_{2i} + \cdots + \gamma_k X_{ki} + u_i$$
\item Extraer el $R^2$ de esta regresión auxiliar
\end{enumerate}
\end{frame}

\begin{frame}{Test de Breusch-Pagan: estadístico y decisión}
\textbf{Procedimiento (continuación):}
\begin{enumerate}\setcounter{enumi}{4}
\item Calcular el estadístico de prueba:
  $$BP = n \cdot R^2 \sim \chi^2_{k}$$
  donde $n$ = tamaño de muestra, $k$ = número de variables explicativas (sin contar la constante)
\item Si $BP > \chi^2_{k, \alpha}$ (o si $p < \alpha$), rechazamos $H_0$ $\Rightarrow$ hay heterocedasticidad
\end{enumerate}
\end{frame}

\begin{frame}[fragile]{Breusch-Pagan en R}
\begin{lstlisting}
library(lmtest)

# Estimar modelo
modelo <- lm(MOD_INGLES_PUNT ~ ESTU_ESTRATO + PRIVADA +
             ESTU_GENERO + PUNT_INGLES_S11, data = ni_data)

# Test de Breusch-Pagan
bptest(modelo)
\end{lstlisting}
\end{frame}

\begin{frame}[fragile]{Breusch-Pagan en R: interpretación}
\begin{exampleblock}{Ejemplo de output}
\small
\begin{verbatim}
    studentized Breusch-Pagan test

data:  modelo
BP = 38.472, df = 4, p-value = 8.741e-08
\end{verbatim}
\end{exampleblock}

\textbf{Interpretación:}
\begin{itemize}
\item $BP = 38.472$ con $df = 4$ (4 variables explicativas)
\item $p < 0.001$ $\Rightarrow$ rechazamos $H_0$
\item \textbf{Conclusión}: hay evidencia de heterocedasticidad
\end{itemize}
\end{frame}

\begin{frame}{Variantes del test de Breusch-Pagan}
\begin{block}{Breusch-Pagan estándar}
Asume que los errores al cuadrado se distribuyen $\chi^2$. Versión más común.
\end{block}

\begin{block}{Breusch-Pagan estudentizado (por defecto en R)}
Usa residuos estudentizados en lugar de residuos crudos. Más robusto a no normalidad.
$BP = n \cdot R^2 \sim \chi^2_{k}$
\end{block}

\begin{block}{Test de White}
Variante más general: incluye cuadrados y productos cruzados de las $X$.
{\small $e_i^2 = \gamma_0 + \gamma_1 X_{1i} + \gamma_2 X_{2i} + \gamma_3 X_{1i}^2 + \gamma_4 X_{2i}^2 + \gamma_5 X_{1i}X_{2i} + u_i$}

Detecta formas más complejas de heterocedasticidad, pero usa más grados de libertad.
\end{block}

\vspace{0.1cm}
\small
\textbf{Regla práctica}: comenzar con Breusch-Pagan. Si hay muchas variables, usar White.
\end{frame}

\begin{frame}{Visualización de heterocedasticidad}
\begin{columns}
\begin{column}{0.48\textwidth}
\textbf{Homocedasticidad (bueno)}
\begin{center}
\begin{tikzpicture}
\begin{axis}[
  width=6cm, height=5cm,
  xlabel={$\hat{Y}$},
  ylabel={Residuos $e_i$},
  grid=major,
  xmin=120, xmax=180,
  ymin=-30, ymax=30
]
\addplot[only marks, mark=*, mark size=1pt, opacity=0.4]
  coordinates {(130,10) (135,-12) (140,8) (145,-10) (150,15) (155,-8) (160,12) (165,-15) (170,10) (140,-5) (150,-10) (160,5)};
\addplot[red, thick, domain=120:180] {0};
\end{axis}
\end{tikzpicture}
\end{center}
\small Dispersión constante alrededor de cero.
\end{column}

\begin{column}{0.48\textwidth}
\textbf{Heterocedasticidad (malo)}
\begin{center}
\begin{tikzpicture}
\begin{axis}[
  width=6cm, height=5cm,
  xlabel={$\hat{Y}$},
  ylabel={Residuos $e_i$},
  grid=major,
  xmin=120, xmax=180,
  ymin=-40, ymax=40
]
\addplot[only marks, mark=*, mark size=1pt, opacity=0.4]
  coordinates {(130,5) (135,-8) (140,10) (145,-12) (150,15) (155,-20) (160,25) (165,-30) (170,35) (140,-5) (150,-10) (160,20)};
\addplot[red, thick, domain=120:180] {0};
\draw[javeriana, dashed, thick] (axis cs:120,-15) -- (axis cs:180,-35);
\draw[javeriana, dashed, thick] (axis cs:120,15) -- (axis cs:180,35);
\end{axis}
\end{tikzpicture}
\end{center}
\small Embudo: dispersión crece con $\hat{Y}$.
\end{column}
\end{columns}
\end{frame}

% ============================================================================
\section{Autocorrelación y Durbin-Watson}
% ============================================================================

\begin{frame}{El supuesto de no autocorrelación}
\begin{block}{Definición}
\textbf{No autocorrelación}: los errores de diferentes observaciones no están correlacionados:
$\text{Cor}(\varepsilon_i, \varepsilon_j) = 0$ para todo $i \neq j$.
\textbf{Autocorrelación}: $\text{Cor}(\varepsilon_i, \varepsilon_{i-1}) \neq 0$.
\end{block}

\vspace{0.1cm}
\begin{columns}
\begin{column}{0.48\textwidth}
\begin{exampleblock}{¿Cuándo es común?}
\begin{itemize}
\item \textbf{Series de tiempo}: observaciones consecutivas (PIB mensual, ventas diarias)
\item \textbf{Datos espaciales}: observaciones vecinas
\item Menos común en corte transversal
\end{itemize}
\end{exampleblock}
\end{column}

\begin{column}{0.48\textwidth}
\begin{alertblock}{Consecuencias}
Similar a heterocedasticidad:
\begin{itemize}
\item $\hat{\beta}$ sigue insesgado
\item Pero ineficiente
\item Errores estándar incorrectos
\item IC y pruebas inválidas
\end{itemize}
\end{alertblock}
\end{column}
\end{columns}
\end{frame}

\begin{frame}{Test de Durbin-Watson}
\begin{block}{Hipótesis}
$H_0$: No hay autocorrelación de primer orden \qquad $H_1$: Hay autocorrelación de primer orden
\end{block}

\textbf{Estadístico de Durbin-Watson:}
$$DW = \frac{\sum_{i=2}^{n}(e_i - e_{i-1})^2}{\sum_{i=1}^{n}e_i^2}$$

\vspace{0.15cm}
\textbf{Interpretación} (rango: $0 \leq DW \leq 4$):
\begin{itemize}
\item $DW \approx 2$: no autocorrelación
\item $DW < 2$: autocorrelación positiva (errores consecutivos con mismo signo)
\item $DW > 2$: autocorrelación negativa (errores consecutivos con signos opuestos)
\end{itemize}

\vspace{0.1cm}
\small
\textbf{Regla práctica}: si $1.5 < DW < 2.5$, probablemente no hay problema.
\end{frame}

\begin{frame}[fragile]{Durbin-Watson en R}
\begin{lstlisting}
library(lmtest)

# Test de Durbin-Watson
dwtest(modelo)
\end{lstlisting}

\vspace{0.2cm}
\begin{exampleblock}{Ejemplo de output}
\small
\begin{verbatim}
    Durbin-Watson test

data:  modelo
DW = 1.9823, p-value = 0.4521
alternative hypothesis: true autocorrelation is greater than 0
\end{verbatim}
\end{exampleblock}
\end{frame}

\begin{frame}{Durbin-Watson: interpretación}
\textbf{Interpretación del ejemplo:}
\begin{itemize}
\item $DW = 1.9823 \approx 2$ $\Rightarrow$ no hay autocorrelación
\item $p = 0.4521 > 0.05$ $\Rightarrow$ no rechazamos $H_0$
\item \textbf{Conclusión}: no hay evidencia de autocorrelación de primer orden
\end{itemize}

\vspace{0.3cm}
\begin{alertblock}{Nota}
En datos de corte transversal (como ICFES), la autocorrelación es rara. El test es más relevante para series de tiempo.
\end{alertblock}
\end{frame}

\begin{frame}{Visualización de autocorrelación}
\begin{columns}
\begin{column}{0.48\textwidth}
\textbf{Sin autocorrelación}
\begin{center}
\begin{tikzpicture}
\begin{axis}[
  width=6cm, height=4.5cm,
  xlabel={Tiempo/Orden},
  ylabel={Residuos},
  grid=major,
  xmin=0, xmax=20,
  ymin=-3, ymax=3
]
\addplot[mark=*, mark size=1.5pt, javerianalight]
  coordinates {(1,1.2) (2,-0.5) (3,0.8) (4,-1.5) (5,0.3) (6,1.0) (7,-0.8) (8,0.5) (9,-1.2) (10,0.7) (11,-0.3) (12,1.5) (13,-0.9) (14,0.4) (15,-1.1)};
\end{axis}
\end{tikzpicture}
\end{center}
\small Patrón aleatorio, sin tendencias.
\end{column}

\begin{column}{0.48\textwidth}
\textbf{Con autocorrelación positiva}
\begin{center}
\begin{tikzpicture}
\begin{axis}[
  width=6cm, height=4.5cm,
  xlabel={Tiempo/Orden},
  ylabel={Residuos},
  grid=major,
  xmin=0, xmax=20,
  ymin=-3, ymax=3
]
\addplot[mark=*, mark size=1.5pt, red]
  coordinates {(1,-1.5) (2,-1.3) (3,-1.0) (4,-0.8) (5,0.2) (6,0.5) (7,0.8) (8,1.2) (9,1.5) (10,1.3) (11,0.8) (12,0.3) (13,-0.5) (14,-1.0) (15,-1.4)};
\end{axis}
\end{tikzpicture}
\end{center}
\small Rachas: positivos seguidos de positivos.
\end{column}
\end{columns}

\vspace{0.3cm}
\small
Con autocorrelación positiva, el DW será significativamente menor que 2. En el ejemplo de la derecha, podríamos tener $DW \approx 0.5$.
\end{frame}

% ============================================================================
\section{Observaciones influyentes: Cook's D}
% ============================================================================

\begin{frame}{Conceptos: leverage y residuos}
\begin{columns}
\begin{column}{0.55\textwidth}
\begin{block}{Leverage (apalancamiento)}
Mide qué tan extremo es el valor de $X_i$ respecto al centro de los datos.
$$h_{ii} = \text{diagonal de } H = X(X'X)^{-1}X'$$
\begin{itemize}
\item $0 \leq h_{ii} \leq 1$, \; $\sum h_{ii} = k+1$
\item Promedio: $\bar{h} = (k+1)/n$
\item Alto leverage: $h_{ii} > 2\bar{h}$
\end{itemize}
\end{block}
\end{column}

\begin{column}{0.42\textwidth}
\begin{block}{Residuo}
Mide qué tan lejos está $Y_i$ de $\hat{Y}_i$:
$$e_i = Y_i - \hat{Y}_i$$
Un residuo grande indica que el modelo no ajusta bien esa observación.
\end{block}
\end{column}
\end{columns}
\end{frame}

\begin{frame}{Distancia de Cook}
\begin{block}{Definición}
Mide cuánto cambiarían TODOS los coeficientes $\hat{\beta}$ si elimináramos la observación $i$.
$$D_i = \frac{e_i^2}{p \cdot MSE} \cdot \frac{h_{ii}}{(1-h_{ii})^2}$$
donde $e_i$ = residuo, $p$ = número de parámetros, $MSE$ = error cuadrático medio, $h_{ii}$ = leverage.
\end{block}

\vspace{0.15cm}
\textbf{Interpretaci\'on:}
\begin{itemize}
\item $D_i$ combina residuo grande y leverage alto
\item Observaci\'on con alto $D_i$ es \textbf{influyente}: afecta mucho los coeficientes
\item Regla de pulgar: $D_i > \frac{4}{n}$ merece investigaci\'on
\item Valores $D_i > 1$ son muy influyentes
\end{itemize}

\smallskip
\small Imaginen un estudiante que falt\'o a todos los parciales pero sac\'o 5.0 en el final --- ese caso distorsiona el modelo de predicci\'on de notas. Eso es una observaci\'on influyente.
\end{frame}

\begin{frame}{Tipos de observaciones}
\begin{center}
\begin{tikzpicture}[scale=0.85]
\begin{axis}[
  width=10cm, height=5.5cm,
  xlabel={Leverage ($h_{ii}$)},
  ylabel={Residuo estandarizado},
  grid=major,
  xmin=0, xmax=0.25,
  ymin=-3, ymax=3,
  legend pos=north east
]

% Observaciones normales
\addplot[only marks, mark=*, mark size=2pt, color=javerianalight, opacity=0.6]
  coordinates {(0.02,0.3) (0.03,-0.5) (0.025,0.8) (0.035,-0.4) (0.04,0.6) (0.03,-0.7)};
\addlegendentry{Normales}

% Outlier
\addplot[only marks, mark=*, mark size=3pt, color=orange]
  coordinates {(0.03, 2.5)};
\addlegendentry{Outlier}

% Alto leverage
\addplot[only marks, mark=*, mark size=3pt, color=blue]
  coordinates {(0.20, 0.3)};
\addlegendentry{Alto leverage}

% Influyente
\addplot[only marks, mark=*, mark size=3pt, color=red]
  coordinates {(0.18, 2.3)};
\addlegendentry{Influyente}

% Lineas de referencia
\addplot[dashed, black, domain=0:0.25] {0};
\end{axis}
\end{tikzpicture}
\end{center}

\vspace{0.2cm}
\small
\begin{itemize}
\item \textcolor{javerianalight}{Normales}: bajo leverage, bajo residuo $\Rightarrow$ $D_i$ bajo
\item \textcolor{orange}{Outlier}: bajo leverage, alto residuo $\Rightarrow$ $D_i$ moderado
\item \textcolor{blue}{Alto leverage}: alto leverage, bajo residuo $\Rightarrow$ $D_i$ bajo (no influyente)
\item \textcolor{red}{Influyente}: alto leverage Y alto residuo $\Rightarrow$ $D_i$ alto
\end{itemize}
\end{frame}

\begin{frame}[fragile]{Calculando la distancia de Cook en R}
\begin{lstlisting}
# Extraer distancia de Cook para cada observacion
cook_d <- cooks.distance(modelo)

# Ver las 5 observaciones mas influyentes
head(sort(cook_d, decreasing = TRUE), 5)

# Identificar observaciones con D > 4/n
n <- nrow(ni_data)
umbral <- 4 / n
influyentes <- which(cook_d > umbral)

# Ver cuantas y cuales observaciones son influyentes
length(influyentes)
ni_data[influyentes, ]
\end{lstlisting}
\end{frame}

\begin{frame}[fragile]{Graficando la distancia de Cook}
\begin{lstlisting}
# Graficar distancia de Cook
plot(cook_d, type = "h", lwd = 2, col = "steelblue",
     main = "Distancia de Cook",
     ylab = "Cook's distance", xlab = "Observacion")
abline(h = umbral, col = "red", lty = 2, lwd = 2)
text(x = n*0.8, y = umbral*1.5, labels = "Umbral 4/n", col = "red")
\end{lstlisting}

\vspace{0.3cm}
\begin{alertblock}{Interpretación del gráfico}
Las barras que sobrepasan la línea roja ($4/n$) corresponden a observaciones influyentes que merecen investigación.
\end{alertblock}
\end{frame}

\begin{frame}{¿Qué hacer con observaciones influyentes?}
\begin{alertblock}{NO eliminar automáticamente}
Eliminar observaciones sin justificación es manipular los datos. Siempre investigar primero.
\end{alertblock}

\textbf{Pasos a seguir:}
\begin{enumerate}
\item \textbf{Verificar datos}: ¿Es un error de entrada de datos? ¿Valor imposible?
\item \textbf{Contextualizar}: ¿Es un caso especial legítimo?
\item \textbf{Re-estimar sin la observación}: ¿Cambian las conclusiones?
\item \textbf{Reportar sensibilidad}: ``Resultados robustos a exclusión de obs. influyentes'' o ``Las conclusiones cambian al excluir\ldots''
\end{enumerate}
\end{frame}

\begin{frame}{Observaciones influyentes: ejemplo}
\begin{exampleblock}{Ejemplo: error de datos}
Si una observaci\'on tiene MOD\_INGLES\_PUNT = 300 (m\'aximo 100), es claramente un error $\Rightarrow$ eliminar. Es como cuando alguien tiene 200 a\~nos en una base de datos --- obviamente algo sali\'o mal.
\end{exampleblock}

\vspace{0.3cm}
\begin{alertblock}{Ejemplo: caso leg\'itimo}
Un estudiante de estrato 1 que sac\'o 95 en ingl\'es: es at\'ipico pero real, como alguien que aprendi\'o ingl\'es viendo series en Netflix sin subt\'itulos. NO eliminar, pero reportar sensibilidad.
\end{alertblock}
\end{frame}

\begin{frame}[fragile]{Análisis de sensibilidad}
\small
\begin{lstlisting}
# Modelo completo
modelo_completo <- lm(MOD_INGLES_PUNT ~ ESTU_ESTRATO + PRIVADA +
                      ESTU_GENERO + PUNT_INGLES_S11, data = ni_data)

# Identificar observaciones influyentes
cook_d <- cooks.distance(modelo_completo)
influyentes <- which(cook_d > 4 / nrow(ni_data))

# Modelo sin observaciones influyentes
ni_sin_inf <- ni_data[-influyentes, ]
modelo_sin_inf <- lm(MOD_INGLES_PUNT ~ ESTU_ESTRATO + PRIVADA +
                     ESTU_GENERO + PUNT_INGLES_S11,
                     data = ni_sin_inf)

# Comparar coeficientes
library(texreg)
screenreg(list(modelo_completo, modelo_sin_inf),
          custom.model.names = c("Completo", "Sin influyentes"))
\end{lstlisting}
\end{frame}

% ============================================================================
\section{Errores est\'andar robustos (Huber-White)}
% ============================================================================

\begin{frame}{La soluci\'on a la heterocedasticidad}
\begin{block}{Errores est\'andar robustos (Huber-White)}
Ante heterocedasticidad, podemos:
\begin{itemize}
\item Transformar las variables para reducir heterocedasticidad (ej: log)
\item Usar WLS (Weighted Least Squares) si conocemos la estructura de $\sigma_i^2$
\item O simplemente usar \textbf{errores est\'andar robustos}
\end{itemize}
\end{block}

\textbf{Ventajas de los errores robustos:}
\begin{itemize}
\item Los coeficientes $\hat{\beta}$ NO cambian (mismos valores que OLS est\'andar)
\item Solo cambian los errores est\'andar (se corrigen)
\item Los IC y pruebas $t$ se vuelven v\'alidos
\item No requiere especificar la forma de heterocedasticidad
\end{itemize}
\end{frame}

\begin{frame}{La soluci\'on a la heterocedasticidad (cont.)}
\begin{alertblock}{Regla pr\'actica}
Si hay heterocedasticidad (BP rechaza $H_0$), SIEMPRE reportar errores robustos. Los errores robustos corrigen las estimaciones cuando la varianza no es constante --- como ponerle un filtro de correcci\'on a una foto que sali\'o sobreexpuesta.
\end{alertblock}
\end{frame}

\begin{frame}{Tipos de errores robustos}
\begin{center}
\renewcommand{\arraystretch}{1.2}
\small
\begin{tabular}{lp{7cm}}
\toprule
\textbf{Tipo} & \textbf{Descripci\'on} \\
\midrule
HC0 & Versi\'on original de White. No corrige por grados de libertad. \\
HC1 & Correcci\'on por grados de libertad: $\frac{n}{n-k}$. M\'as com\'un en econometr\'ia. \\
HC2 & Usa leverage: $\frac{e_i^2}{1 - h_{ii}}$. M\'as robusto. \\
HC3 & Versi\'on m\'as conservadora: $\frac{e_i^2}{(1 - h_{ii})^2}$. Recomendado para $n$ peque\~no. \\
\bottomrule
\end{tabular}
\end{center}

\begin{block}{?`Cu\'al usar?}
\begin{itemize}
\item \textbf{HC1}: el más común, recomendado por defecto
\item \textbf{HC3}: si la muestra es pequeña ($n < 250$) o hay alto leverage
\item En la práctica, los resultados suelen ser muy similares
\end{itemize}
\end{block}
\end{frame}

\begin{frame}[fragile]{Errores robustos en R}
\begin{lstlisting}
library(sandwich)  # para vcovHC
library(lmtest)    # para coeftest

# Modelo OLS estandar
modelo <- lm(MOD_INGLES_PUNT ~ ESTU_ESTRATO + PRIVADA +
             ESTU_GENERO + PUNT_INGLES_S11, data = ni_data)

# Tabla con errores estandar clasicos (NO robustos)
summary(modelo)

# Tabla con errores estandar robustos (HC1)
coeftest(modelo, vcov = vcovHC(modelo, type = "HC1"))

# Tambien puedes usar HC3
coeftest(modelo, vcov = vcovHC(modelo, type = "HC3"))
\end{lstlisting}

\vspace{0.15cm}
\small
\textbf{Nota}: \code{vcovHC} calcula la matriz de varianza-covarianza robusta, y \code{coeftest} la usa para recalcular los errores estándar y las pruebas $t$.
\end{frame}

\begin{frame}[fragile]{Comparación: errores estándar clásicos}
\begin{exampleblock}{Errores estándar clásicos (OLS)}
\small
\begin{verbatim}
                    Estimate Std. Error t value Pr(>|t|)
(Intercept)         120.345      3.215  37.434  < 2e-16 ***
ESTU_ESTRATO          2.134      0.421   5.068 4.32e-07 ***
PRIVADA               8.765      1.234   7.103 1.56e-12 ***
ESTU_GENERO          -1.543      0.987  -1.563    0.118
PUNT_INGLES_S11       0.512      0.034  15.059  < 2e-16 ***
\end{verbatim}
\end{exampleblock}
\end{frame}

\begin{frame}[fragile]{Comparación: errores estándar robustos}
\begin{exampleblock}{Errores estándar robustos (HC1)}
\small
\begin{verbatim}
                    Estimate Std. Error t value Pr(>|t|)
(Intercept)         120.345      3.654  32.936  < 2e-16 ***
ESTU_ESTRATO          2.134      0.512   4.168 3.18e-05 ***
PRIVADA               8.765      1.543   5.681 1.52e-08 ***
ESTU_GENERO          -1.543      1.087  -1.419    0.156
PUNT_INGLES_S11       0.512      0.041  12.488  < 2e-16 ***
\end{verbatim}
\end{exampleblock}

\vspace{0.3cm}
Los coeficientes son idénticos, pero los SE aumentaron (especialmente para PRIVADA). Algunas variables pueden perder significancia.
\end{frame}

\begin{frame}[fragile]{Intervalos de confianza robustos}
\begin{lstlisting}
# IC con errores estandar clasicos
confint(modelo)

# IC con errores robustos HC1
coefci(modelo, vcov = vcovHC(modelo, type = "HC1"))
\end{lstlisting}

\vspace{0.15cm}
\begin{columns}
\begin{column}{0.48\textwidth}
\begin{exampleblock}{IC clásicos (95\%)}
\small
\begin{verbatim}
              2.5 %  97.5 %
ESTU_ESTRATO  1.309   2.959
PRIVADA       6.347  11.183
\end{verbatim}
\end{exampleblock}
\end{column}

\begin{column}{0.48\textwidth}
\begin{exampleblock}{IC robustos (95\%)}
\small
\begin{verbatim}
              2.5 %  97.5 %
ESTU_ESTRATO  1.130   3.138
PRIVADA       5.740  11.790
\end{verbatim}
\end{exampleblock}
\end{column}
\end{columns}

\vspace{0.1cm}
\small
Los IC robustos son más amplios (reflejan la mayor incertidumbre por heterocedasticidad).
\end{frame}

\begin{frame}{¿Cuándo usar errores robustos?}
\begin{center}
\renewcommand{\arraystretch}{1.4}
\begin{tabular}{lcc}
\toprule
\textbf{Situación} & \textbf{Acción} & \textbf{Razón} \\
\midrule
Breusch-Pagan rechaza $H_0$ & Usar robustos & Heterocedasticidad detectada \\
Breusch-Pagan no rechaza $H_0$ & Opcional & No hay evidencia de heteroced. \\
Muestra grande ($n > 1000$) & Usar robustos & Casi no cuesta nada \\
Publicación académica & Usar robustos & Estándar en econometría \\
\bottomrule
\end{tabular}
\end{center}

\vspace{0.15cm}
\begin{block}{Recomendación general}
En econometría moderna, es estándar reportar \textbf{siempre} errores robustos, incluso si Breusch-Pagan no rechaza $H_0$.
\end{block}
\end{frame}

\begin{frame}{Limitaciones de los errores robustos}
\begin{alertblock}{Importante}
Los errores robustos NO corrigen sesgos. Si el modelo está mal especificado (variable omitida, no linealidad), los errores robustos no ayudan.
\end{alertblock}
\end{frame}

% ============================================================================
\section{Ejemplo integrado en R}
% ============================================================================

\begin{frame}[fragile]{Ejemplo completo: diagnósticos paso a paso}
\small
\begin{lstlisting}
# 1. Cargar datos
ni_data <- read.csv("icfes_ni.csv")

# 2. Estimar modelo
modelo <- lm(MOD_INGLES_PUNT ~ ESTU_ESTRATO + PRIVADA +
             ESTU_GENERO + PUNT_INGLES_S11, data = ni_data)

# 3. Resumen basico
summary(modelo)

# 4. Graficos diagnosticos
par(mfrow = c(2, 2))
plot(modelo)
par(mfrow = c(1, 1))
# Revisar: patron/embudo, Q-Q, Scale-Location, Cook
\end{lstlisting}
\end{frame}

\begin{frame}[fragile]{Ejemplo completo: tests formales}
\small
\begin{lstlisting}
library(lmtest); library(car)

# 5. Heterocedasticidad (p < 0.05 => heteroced.)
bptest(modelo)

# 6. Autocorrelacion (DW ~ 2 => ok)
dwtest(modelo)

# 7. Multicolinealidad (VIF > 10 => severa)
vif(modelo)

# 8. Observaciones influyentes
cook_d <- cooks.distance(modelo)
influyentes <- which(cook_d > 4 / nrow(ni_data))
length(influyentes)
ni_data[influyentes, ]
\end{lstlisting}
\end{frame}

\begin{frame}[fragile]{Ejemplo completo: corrección con errores robustos}
\small
\begin{lstlisting}
library(sandwich); library(texreg)

# 9. Errores robustos
coeftest(modelo, vcov = vcovHC(modelo, type = "HC1"))

# 10. Comparar clasicos vs robustos
screenreg(list(modelo, modelo),
  override.se = list(NULL, sqrt(diag(vcovHC(modelo, "HC1")))),
  override.pvalues = list(NULL,
    coeftest(modelo, vcovHC(modelo, "HC1"))[, 4]),
  custom.model.names = c("OLS clasico", "OLS robusto"))
\end{lstlisting}
\end{frame}

\begin{frame}[fragile]{Ejemplo completo: análisis de sensibilidad}
\small
\begin{lstlisting}
# 11. Analisis de sensibilidad (sin influyentes)
if (length(influyentes) > 0) {
  modelo_sin_inf <- lm(MOD_INGLES_PUNT ~ ESTU_ESTRATO + PRIVADA +
                       ESTU_GENERO + PUNT_INGLES_S11,
                       data = ni_data[-influyentes, ])
  screenreg(list(modelo, modelo_sin_inf),
            custom.model.names = c("Completo", "Sin influyentes"))
}
\end{lstlisting}
\end{frame}

\begin{frame}{Interpretación del ejemplo}
\textbf{Hallazgos esperados con datos ICFES de NI:}
\begin{enumerate}
\item \textbf{Heterocedasticidad}: BP probablemente rechaza $H_0$. Estratos bajos tienen mayor variabilidad $\Rightarrow$ errores robustos HC1
\item \textbf{Normalidad}: Q-Q con desviación leve en colas. No crítico con $n$ grande (TCL)
\item \textbf{Obs.\ influyentes}: probablemente algunas (5-10). Verificar si son errores o casos legítimos; reportar sensibilidad
\item \textbf{Autocorrelación}: DW $\approx 2$ (no hay problema en corte transversal)
\item \textbf{Multicolinealidad}: VIF probablemente bajo (variables poco correlacionadas)
\end{enumerate}
\end{frame}

% ============================================================================
\section{Resumen}
% ============================================================================

\begin{frame}{Checklist de diagnósticos}
\begin{block}{Protocolo estándar -- después de estimar cualquier modelo OLS:}
\begin{enumerate}
\item Generar \code{plot(modelo)} y revisar los 4 gráficos
\item Realizar \code{bptest(modelo)} para heterocedasticidad
\item Si hay heterocedasticidad, usar \code{coeftest(modelo, vcovHC(...))}
\item Verificar \code{vif(modelo)} para multicolinealidad
\item Identificar obs.\ influyentes con \code{cooks.distance(modelo)}
\item Investigar influyentes y reportar análisis de sensibilidad
\item Reportar resultados con errores robustos si hay heterocedasticidad
\end{enumerate}
\end{block}

\vspace{0.1cm}
\begin{alertblock}{Importante}
Nunca reportar resultados OLS sin verificar supuestos. Los diagnósticos son parte fundamental del análisis.
\end{alertblock}
\end{frame}

\begin{frame}{Tabla resumen: problemas y soluciones}
\footnotesize
\begin{center}
\renewcommand{\arraystretch}{1.2}
\begin{tabular}{llll}
\toprule
\textbf{Problema} & \textbf{Test} & \textbf{Detección} & \textbf{Solución} \\
\midrule
Heterocedasticidad & Breusch-Pagan & Embudo en residuos, & Errores robustos \\
 &  & Scale-Location con pendiente & (HC1 o HC3) \\
\midrule
Autocorrelación & Durbin-Watson & Rachas en residuos, & Errores robustos \\
 &  & $DW \ll 2$ o $DW \gg 2$ & (HAC) \\
\midrule
No normalidad & Q-Q plot & Desviación de línea, colas pesadas & Transformar Y / TCL \\
\midrule
No linealidad & Res. vs Fitted & Patrón curvo, línea roja curvada & Transformar X o Y \\
\midrule
Obs. influyentes & Cook's D & $D_i > 4/n$, puntos extremos & Investigar, sensibilidad \\
\midrule
Multicolinealidad & VIF & $VIF > 10$ & Eliminar var. / PCA \\
\bottomrule
\end{tabular}
\end{center}
\end{frame}

\begin{frame}{Pr\'oxima sesi\'on}
\begin{block}{Sesi\'on 12: Integraci\'on y An\'alisis Aplicado}
Aplicaremos TODO lo aprendido en el curso a un an\'alisis completo con datos ICFES:
\begin{itemize}
\item Pipeline estad\'istico completo
\item An\'alisis de NI: ?`qu\'e determina el puntaje de ingl\'es?
\item Interpretaci\'on de resultados para pol\'itica educativa
\item Retroalimentaci\'on del Problem Set 8
\end{itemize}
\end{block}

\begin{alertblock}{Para la pr\'oxima clase}
\textbf{Traer:}
\begin{itemize}
\item Laptop con R y los paquetes: lmtest, sandwich, car, texreg
\item Datos ICFES de NI (disponibles en Moodle)
\item Dudas sobre diagn\'osticos de regresi\'on
\end{itemize}
\end{alertblock}
\end{frame}

\begin{frame}{}
\begin{center}
\Huge !`Gracias!

\vspace{0.8cm}
\Large Preguntas y discusi\'on

\vspace{1cm}
\normalsize
\textbf{Contacto:}\\
jaime.polanco@javeriana.edu.co\\
Horario de atenci\'on: Mi\'ercoles 14:00-16:00
\end{center}
\end{frame}

\end{document}
